% problem-000166 8 有机反应负氢迁移
\section*{8. 有机反应负氢迁移}
以芳⾹醛和1-苯基⼄-1,2-⼆醇为原料在酸性条件下可以构建环戊烷-2-烯-1-酮衍⽣物：

$
\mathrm{Ph} - \underset {\mathrm{OH}} {\overset {\mathrm{OH}} {\mathrm{C}}} + \mathrm{PhCHO} \xrightarrow [ \mathrm{PhCH} _ {3} , 1 1 0 ^ {\circ} \mathrm{C} ]{p - \mathrm{TsOH(0.5eq)}} \mathrm{Ph} - \underset {\mathrm{Ph}} {\overset {\mathrm{O}} {\mathrm{C}}} - \underset {\mathrm{Ph}} {\overset {\mathrm{Ph}} {\mathrm{C}}}
$

该反应在对甲苯磺酸 $( p { \mathrm { - } } \mathrm { T s O H } )$ 催化下，在甲苯中回流后得到产物（产率70%）。研究表明，苯基环氧⼄烷代替1-苯基⼄-1,2-⼆醇在相同条件下反应也转化为同⼀产物（产率75%）。依据这些信息，回答以下问题：

\subsection*{8-1}
在酸性条件下，1-苯基⼄-1,2-⼆醇和苯基环氧⼄烷均转化为同⼀中间产物，此中间产物随后与苯甲醛反应转化为最终产物2,4,5-三苯基环戊烷-2-烯-1-酮。画出此中间产物的结构式。

\subsection*{8-2}
此反应中苯基环氧⼄烷与苯甲醛的⽐例⾄少应该为多少？

\subsection*{8-3}
以下转换经历了类似的过程，为以下转换提供关键中间体：

（图：）

\subsection*{8-4}
依据8-1和8-2的信息，请完成以下反应式（要求⽴体化学）：

（图：）
